%% ---------------------------------------------------------------------
%% CHAPTER TEMPLATE  (copy to Overleaf/chapters/chNN-slug.tex and fill in)
%% Every chapter of the book follows this order.  Keep the section
%% headings; rename them to fit the topic, but do not drop a stage.
%% ---------------------------------------------------------------------
\chapter{Algorithm Name}
\label{ch:chNN}
% Function names for pseudocode must be unique across the book: prefix them
% with a chapter-specific word.
\SetKwFunction{AlgMain}{AlgorithmName}
\SetKwFunction{AlgHelper}{HelperName}

\begin{objectives}
\begin{itemize}
  \item Explain in one paragraph what the algorithm computes and when it is the right tool.
  \item Trace the algorithm by hand on a small instance.
  \item State the conditions under which it is complete / optimal, and why.
  \item Implement it in Python and test it on a grid or a 2D scene.
  \item Recognise the failure modes that matter for a drone swarm.
\end{itemize}
\end{objectives}

% ---------------------------------------------------------------------
\section{Why this matters for a drone swarm}
\label{sec:chNN-motivation}
% 1/2 to 1 page.  A concrete drone scenario in which this algorithm is needed.
% Where it lives in the four-layer hybrid architecture (\cref{ch:ch24}).
% End with a short roadmap of the chapter.

% ---------------------------------------------------------------------
\section{The idea in plain words}
\label{sec:chNN-intuition}
% Intuition before formalism.  A picture (\inputfigure{chNN/idea}) and
\begin{keyidea}
One or two sentences that a reader should remember for the rest of their life.
\end{keyidea}

% ---------------------------------------------------------------------
\section{Problem statement and notation}
\label{sec:chNN-problem}
% Formal definitions.  Use the macros of searchbook.sty (\gcost, \hcost, \pos, ...).
\begin{definition}[Name]
\label{def:chNN-name}
...
\end{definition}

% ---------------------------------------------------------------------
\section{The algorithm}
\label{sec:chNN-algorithm}
% Pseudocode + line-by-line walkthrough (refer to lines with \ref{alg:...:line}).
\begin{algorithm}[htb]
\caption{Algorithm name}
\label{alg:chNN-name}
\KwIn{...}
\KwOut{...}
\Fn{\AlgMain{$s$, $g$}}{
  initialise \Open with $s$\;
  \While{\Open is not empty}{
    $n \gets$ node in \Open with the smallest $\fcost$\label{alg:chNN-name:pop}\;
    \If{$n = g$}{\KwRet path\;}
    \ForEach{$n' \in \Succ(n)$}{ ... }
  }
  \KwRet failure\;
}
\end{algorithm}

% ---------------------------------------------------------------------
\section{A worked example}
\label{sec:chNN-example}
% A small, fully worked instance with a figure and a trace table.
\begin{example}[Descriptive title]
\label{ex:chNN-name}
...
\end{example}
\begin{figure}[tb]
  \centering
  \inputfigure{chNN/example}
  \caption{What the reader should notice in this figure.}
  \label{fig:chNN-example}
\end{figure}
\begin{table}[tb]
  \centering\small
  \caption{Trace of \cref{alg:chNN-name} on \cref{ex:chNN-name}.}
  \label{tab:chNN-trace}
  \begin{tabular}{@{}rlll@{}}
  \toprule
  Step & Node expanded & \Open after the step & Comment\\
  \midrule
  1 & ... & ... & ...\\
  \bottomrule
  \end{tabular}
\end{table}

% ---------------------------------------------------------------------
\section{Properties: correctness, optimality, complexity}
\label{sec:chNN-properties}
\begin{theorem}[Name]
\label{thm:chNN-name}
...
\end{theorem}
\begin{proof}
... (a complete proof for the essential results; a proof sketch with a citation otherwise)
\end{proof}

% ---------------------------------------------------------------------
\section{Variants and extensions}
\label{sec:chNN-variants}

% ---------------------------------------------------------------------
\section{Implementation notes}
\label{sec:chNN-implementation}
% Data structures, tie-breaking, numerical issues, then a listing excerpted
% from code/chNN_slug.py (ASCII only, at most 45 lines per listing).
\begin{lstlisting}[caption={Core of the implementation (from \code{code/chNN\_slug.py}).},label={lst:chNN-core}]
def algorithm(start, goal):
    ...
\end{lstlisting}

\begin{pitfall}
A common mistake, why it happens, and how to avoid it.
\end{pitfall}

% ---------------------------------------------------------------------
\section{Where this fits in the drone system}
\label{sec:chNN-drone}
\begin{dronebox}
How the algorithm is used in the hybrid architecture of \cref{ch:ch24}, what it
receives from and delivers to the other layers, and which week of the study plan
(\cref{ch:appA}) covers it.
\end{dronebox}

% ---------------------------------------------------------------------
\section{Summary}
\begin{summary}
\begin{itemize}
  \item ...
\end{itemize}
\end{summary}

\paragraph{Further reading.}
Original source \cite{key}; textbook treatment \cite{key}.

% ---------------------------------------------------------------------
\section{Exercises}
\label{sec:chNN-exercises}
\begin{exercise}[\difficulty{1}]
\label{exr:chNN-slug}
...
\end{exercise}
\begin{exercise}[\difficulty{3} Coding]
\label{exr:chNN-coding}
The coding exercise of the corresponding week in the study plan.
\end{exercise}
